\documentclass[10pt]{beamer}
\usepackage{../../../resources/themes/algolymp-dense}

\title{Divizibilitate}
\subtitle{Ciurul lui Eratostene $\cdot$ Descompunere în factori primi $\cdot$ Suma și numărul divizorilor}
\author{Andrei Mihai Ciobanu}
\institute{Algolymp Summer School}
\date{}
\titlegraphic{\includegraphics[width=3cm]{../../../resources/algolymp-logo.png}}
\setbeamertemplate{frame footer}{Algolymp Summer School}

% ============================================================
\begin{document}

\begin{frame}\titlepage\end{frame}

\begin{frame}{Cuprins}\tableofcontents\end{frame}

% ============================================================
\section{Divizibilitate}

\begin{frame}{Divizibilitate}
  \begin{block}{Definiție}
    Spunem că $a$ \textbf{îl divide} pe $b$ (scriem $a \mid b$) dacă există un număr întreg $c$
    astfel încât $b = a \cdot c$.
  \end{block}
  \begin{itemize}
    \item $a$ se numește \textbf{divizor} al lui $b$
    \item $b$ se numește \textbf{multiplu} al lui $a$
  \end{itemize}
  \smallskip
  Exemple:
  \begin{itemize}
    \item $3 \mid 12$ deoarece $12 = 3 \cdot 4$
    \item $5 \mid 30$ deoarece $30 = 5 \cdot 6$
    \item $4 \nmid 10$ deoarece nu există $c$ întreg cu $10 = 4 \cdot c$
  \end{itemize}
\end{frame}

\begin{frame}{Divizori proprii și improprii}
  \begin{block}{Definiție}
    \textbf{Divizorii proprii} ai lui $n$ sunt toți divizorii săi \textbf{diferiți de $1$ și de $n$}.\\
    \textbf{Divizorii improprii} ai lui $n$ sunt $1$ și $n$ însuși.
  \end{block}
  Exemplu: $n = 12$, divizorii sunt $1, 2, 3, 4, 6, 12$
  \begin{itemize}
    \item Proprii: $2, 3, 4, 6$
    \item Improprii: $1, 12$
  \end{itemize}
\end{frame}

\begin{frame}{Găsirea divizorilor \textemdash{} varianta $O(n)$}
  \begin{block}{Idee}
    Verificăm fiecare $d \in \{1, 2, \ldots, n\}$ dacă $d \mid n$.
  \end{block}
  Exemplu: $n = 12$ $\Rightarrow$ divizori găsiți: $1, 2, 3, 4, 6, 12$
  \begin{itemize}
    \item Corect, dar ineficient pentru $n$ mare ($n > 10^7$)
    \item Putem face mai bine?
  \end{itemize}
  \beginnote{$O(f(n))$ = notație pentru complexitate; descrie câte operații face algoritmul raportat la $n$}
\end{frame}

\begin{frame}{Găsirea divizorilor \textemdash{} reducere la $\lfloor n/2 \rfloor$ iterații}
  Exemplu: $n = 12$, divizorii sunt $1, 2, 3, 4, 6, 12$.
  Cel mai mare divizor propriu este $6 = 12/2$. De ce nu poate fi mai mare?
  \begin{block}{Demonstrație}
    Dacă $d \mid n$, atunci $n = d \cdot k$ cu $k \in \mathbb{N}^*$.\\
    Dacă $d \neq n$, atunci $k \geq 2$, deci $d = \dfrac{n}{k} \leq \dfrac{n}{2}$.
  \end{block}
  \begin{itemize}
    \item Căutăm în $\left\{1, 2, \ldots, \left\lfloor\dfrac{n}{2}\right\rfloor\right\}$ și adăugăm $n$ separat
    \item De 2 ori mai rapid, dar încă ineficient pentru $n$ mare
  \end{itemize}
\end{frame}

\begin{frame}{Găsirea divizorilor \textemdash{} optimizare la $O(\sqrt{n})$}
  \begin{block}{Observație cheie}
    Dacă $d \mid n$, atunci și $\frac{n}{d} \mid n$.\\
    Divizorii vin în \textbf{perechi} $(d,\ n/d)$, cu $d \leq \sqrt{n} \leq n/d$.
  \end{block}
  \begin{itemize}
    \item Căutăm doar $d \in \{1, 2, \ldots, \lfloor\sqrt{n}\rfloor\}$
    \item Pentru fiecare $d$ găsit, adăugăm și perechea $n/d$
    \item \alert{Atenție:} dacă $d = \sqrt{n}$ (pătrat perfect), adăugăm $d$ o singură dată
  \end{itemize}
  \smallskip
  Exemplu: $n = 36$, $\sqrt{36} = 6$\\
  $d=1\to(1,36)$, $d=2\to(2,18)$, $d=3\to(3,12)$, $d=4\to(4,9)$, $d=6\to(6)$
\end{frame}

\begin{frame}[fragile]{Găsirea divizorilor \textemdash{} implementare $O(\sqrt{n})$}
  \begin{lstlisting}
void find_divisors(int n) {
    for (int d = 1; d * d <= n; d++) {
        if (n % d == 0) {
            cout << d << " ";
            if (d != n / d) {
                cout << n / d << " ";
            }
        }
    }
}
  \end{lstlisting}
  \begin{itemize}
    \item Divizorii nu sunt găsiți în ordine crescătoare
    \item \alert{Overflow:} \lstinline|d * d| poate depăși \lstinline|INT_MAX| $(2^{31}-1)$ când \lstinline|n| se apropie de \lstinline|INT_MAX| $\rightarrow$ fix rapid: \lstinline|(long long)d * d <= n|
    \item Timp: $O(\sqrt{n})$, spațiu: $O(1)$
  \end{itemize}
  \beginnote{\lstinline|void find_divisors(int n)| este o \textbf{funcție} C++ care primește $n$ și afișează divizorii}
\end{frame}

% ============================================================
\section{Numărul și suma divizorilor}

\begin{frame}{Numărul și suma divizorilor în $O(\sqrt{n})$}
  Folosind aceeași abordare, calculăm numărul divizorilor ($d(n)$ sau $\tau(n)$) și suma divizorilor ($\sigma(n)$):
  \begin{itemize}
    \item Pentru fiecare pereche $(d,\ n/d)$ cu $d < n/d$: adăugăm \textbf{2} la $\tau$ și $d + n/d$ la $\sigma$
    \item Dacă $d = n/d$ (pătrat perfect): adăugăm \textbf{1} la $\tau$ și $d$ la $\sigma$
  \end{itemize}
  \begin{block}{Observație \textemdash{} pătrate perfecte}
    $n$ este pătrat perfect $\Leftrightarrow$ $\tau(n)$ este \textbf{impar}.\\
    Toate celelalte numere au număr par de divizori.
  \end{block}
\end{frame}

\begin{frame}[fragile]{Numărul și suma divizorilor \textemdash{} implementare}
  \begin{lstlisting}
pair<int, long long> nrdiv_sumdiv(int n) {
    int nrdiv = 0;
    long long sumdiv = 0;
    for (int d = 1; d * d <= n; d++) {
        if (n % d == 0) {
            if (d == n / d) {
                nrdiv += 1;
                sumdiv += d;
            } else {
                nrdiv += 2;
                sumdiv += d + n / d;
            }
        }
    }
    return { nrdiv, sumdiv };
}
  \end{lstlisting}
  Exemplu: $n=12 \Rightarrow \tau(n)=6,\ \sigma(n)=28$
  \beginnote{\lstinline|pair<int, long long>| = tip STL care stochează 2 valori; \lstinline|.first| și \lstinline|.second| le accesează}
\end{frame}

% ============================================================
\section{Numere prime și verificarea primalității}

\begin{frame}{Numere prime și compuse}
  \begin{block}{Număr prim}
    Un număr natural $p \geq 2$ este \textbf{prim} dacă are exact \textbf{2} divizori: $1$ și $p$.
  \end{block}
  \begin{block}{Număr compus}
    Un număr natural $n \geq 2$ este \textbf{compus} dacă are \textbf{mai mult de 2} divizori.
  \end{block}
  \begin{itemize}
    \item Prime: $2, 3, 5, 7, 11, 13, 17, 19, \ldots$
    \item Compuse: $4, 6, 8, 9, 10, 12, \ldots$
    \item $0$ și $1$ nu sunt \textbf{nici prime, nici compuse}
  \end{itemize}
\end{frame}

\begin{frame}{Verificarea primalității}
  \begin{block}{Idee}
    $n$ este prim $\Leftrightarrow$ nu are niciun divizor $d$ cu $1 < d < n$.
  \end{block}
  Aplicăm din nou optimizarea cu $\sqrt{n}$:
  \begin{itemize}
    \item Dacă $n$ are măcar un divizor propriu, atunci are măcar un divizor propriu $\leq \sqrt{n}$
    \item Căutăm în $\{2, 3, \ldots, \lfloor\sqrt{n}\rfloor\}$
    \item Dacă nu găsim niciunul, $n$ este prim
  \end{itemize}
  \begin{alertblock}{Cazuri speciale}
    $n < 2$ $\rightarrow$ nici prim, nici compus \quad $|$ \quad $n = 2$ $\rightarrow$ prim
  \end{alertblock}
\end{frame}

\begin{frame}[fragile]{Verificarea primalității \textemdash{} implementare}
  \begin{lstlisting}
bool is_prime(int n) {
    if (n < 2) {
        return false;
    }
    for (int d = 2; d * d <= n; d++) {
        if (n % d == 0) {
            return false;
        }
    }
    return true;
}
  \end{lstlisting}
  \begin{itemize}
    \item Timp: $O(\sqrt{n})$, spațiu: $O(1)$
    \item Eficient pentru un singur număr; ineficient dacă testăm toate numerele până la $n$
  \end{itemize}
\end{frame}

% ============================================================
\section{Descompunerea în factori primi}

\begin{frame}{Factori primi}
  \begin{block}{Observație cheie}
    Cel mai mic divizor al lui $n$ mai mare ca $1$ este \textbf{întotdeauna prim}.
  \end{block}
  \textbf{De ce?} Fie $p$ cel mai mic divizor $> 1$ al lui $n$. Dacă $p$ ar fi compus,
  ar avea un divizor $q$ cu $1 < q < p$. Dar atunci $q \mid p$ și $p \mid n$, deci $q \mid n$
  \textemdash{} contradicție, deoarece $p$ nu ar mai fi cel mai mic divizor al lui $n$.
  \begin{block}{Teorema fundamentală a aritmeticii}
    Orice $n \geq 2$ se scrie \textbf{unic} ca produs de puteri de prime:
    \[
      n = p_1^{a_1} \cdot p_2^{a_2} \cdots p_k^{a_k}, \quad p_1 < p_2 < \cdots < p_k
    \]
  \end{block}
\end{frame}

\begin{frame}{Algoritm de descompunere}
  \begin{block}{Idee}
    Împărțim repetat $n$ la cel mai mic factor $p \leq \sqrt{n}$.\\
    La final, dacă $n > 1$, restul $n$ este un factor prim cu exponentul $1$.
  \end{block}
  Exemplu: $n = 360$
  \begin{itemize}
    \item $360 \div 2 = 180 \div 2 = 90 \div 2 = 45 \quad\Rightarrow\quad 2^3$
    \item $45 \div 3 = 15 \div 3 = 5 \quad\Rightarrow\quad 3^2$
    \item $5 > 1 \quad\Rightarrow\quad 5^1$
    \item Rezultat: $360 = 2^3 \cdot 3^2 \cdot 5$
  \end{itemize}
\end{frame}

\begin{frame}[fragile]{Descompunere în factori primi \textemdash{} implementare}
  \begin{lstlisting}
vector<pair<int, int>> factorize(int n) {
    vector<pair<int, int>> factors; // (factor prim, exponent)
    for (int p = 2; p * p <= n; p++) {
        if (n % p == 0) {
            int e = 0;
            while (n % p == 0) {
                e++;
                n /= p;
            }
            factors.push_back({ p, e });
        }
    }
    if (n > 1) {
        factors.push_back({ n, 1 });
    }
    return factors;
}
  \end{lstlisting}
  Timp: $O(\sqrt{n})$, spațiu: $O(\log n)$ \textemdash{} cel mult $\log_2 n$ factori primi distincți
  \beginnote{$\log_2 n$ = de câte ori împarți $n$ la $2$ până ajungi la $1$; $\log_2(10^6)\approx 20$. \lstinline|vector| = listă dinamică (STL)}
\end{frame}

\begin{frame}{De la descompunere la $\tau(n)$ și $\sigma(n)$}
  Fie $n = p_1^{a_1} \cdot p_2^{a_2} \cdots p_k^{a_k}$. Orice divizor al lui $n$ are forma
  $p_1^{b_1} \cdot p_2^{b_2} \cdots p_k^{b_k}$ cu $0 \leq b_i \leq a_i$.
  \begin{block}{Numărul de divizori}
    Fiecare $b_i$ poate lua $a_i + 1$ valori, deci:
    \[ \tau(n) = (a_1+1)(a_2+1)\cdots(a_k+1) \]
  \end{block}
  \begin{block}{Suma divizorilor}
    Înmulțind contribuția fiecărui factor prim:
    \[ \sigma(n) = \prod_{i=1}^{k}(1 + p_i + p_i^2 + \cdots + p_i^{a_i}) = \prod_{i=1}^{k}\frac{p_i^{a_i+1}-1}{p_i-1} \]
  \end{block}
  Exemplu: $360 = 2^3 \cdot 3^2 \cdot 5^1$\\
  $\tau(360) = 4 \cdot 3 \cdot 2 = 24$ \quad
  $\sigma(360) = (1{+}2{+}4{+}8)(1{+}3{+}9)(1{+}5) = 1170$
\end{frame}

% ============================================================
\section{Ciurul lui Eratostene}

\begin{frame}{Motivație \textemdash{} de ce avem nevoie de ciur?}
  \begin{itemize}
    \item \texttt{is\_prime(n)} în $O(\sqrt{n})$ e bun pentru \textbf{un singur număr}
    \item Dacă vrem \textbf{toate} primele $\leq n$, aplicând \texttt{is\_prime} de $n$ ori obținem $O(n\sqrt{n})$
    \item Putem face mult mai bine: $O(n \log n)$ sau chiar $O(n \log \log n)$
  \end{itemize}
\end{frame}

\begin{frame}{Ciurul lui Eratostene \textemdash{} idee}
  \begin{block}{Principiu}
    Pornim cu toate numerele $\geq 2$ marcate ca \textbf{prime}.\\
    Pentru fiecare prim $p$ găsit, marcăm toți multiplii $m \geq 2p$ ai lui $p$ ca \textbf{compuși}.
  \end{block}
  \begin{enumerate}
    \item Inițializăm \texttt{prime[i] = true} pentru $i \geq 2$
    \item Pentru $i = 2, 3, \ldots, \lfloor\sqrt{n}\rfloor$: dacă \texttt{prime[i]}, marcăm $2i,\ 3i,\ 4i,\ \ldots$ cu \texttt{false}
    \item \texttt{prime[k] == true} $\Leftrightarrow$ $k$ este prim
  \end{enumerate}
  \begin{alertblock}{Optimizare}
    Putem porni marcarea de la $i^2$ în loc de $2i$, deoarece multiplii $i \cdot k$ cu $k < i$ au fost deja marcați la pasul $k$.
  \end{alertblock}
\end{frame}

\begin{frame}[fragile]{Ciurul lui Eratostene \textemdash{} implementare}
  \begin{lstlisting}
const int nmax = 1000000;
bool prime[nmax + 1];

void sieve() {
    for (int i = 2; i <= nmax; i++) {
        prime[i] = true;
    }
    for (int i = 2; i * i <= nmax; i++) {
        if (prime[i]) {
            for (int j = 2 * i; j <= nmax; j += i) {
                prime[j] = false;
            }
        }
    }
}
  \end{lstlisting}
  Timp: $O(n \log n)$, spațiu: $O(n)$
\end{frame}

\begin{frame}[fragile]{Extensie \textemdash{} ciur pentru $\tau(n)$ și $\sigma(n)$}
  Modificăm ciurul pentru a calcula $\tau$ și $\sigma$ pentru \textbf{toate} numerele până la $n$:
  \begin{lstlisting}
const int nmax = 1000000;
int nrdiv[nmax + 1];
long long sumdiv[nmax + 1];

void sieve_divisors() {
    for (int d = 1; d <= nmax; d++) {
        for (int m = d; m <= nmax; m += d) {
            nrdiv[m]++;
            sumdiv[m] += d;
        }
    }
}
  \end{lstlisting}
  Timp: $O(n \log n)$, spațiu: $O(n)$
\end{frame}

% ============================================================
\section{Recapitulare}

\begin{frame}{Recapitulare}
  \begin{center}
  \begin{tabular}{lll}
    \toprule
    Algoritm & Timp & Spațiu \\
    \midrule
    Găsire divizori                   & $O(\sqrt{n})$ & $O(1)$ \\
    $\tau(n)$, $\sigma(n)$ (un număr) & $O(\sqrt{n})$ & $O(1)$ \\
    Verificare primalitate            & $O(\sqrt{n})$ & $O(1)$ \\
    Descompunere în factori primi     & $O(\sqrt{n})$ & $O(\log n)$ \\
    Ciurul lui Eratostene             & $O(n \log n)$ & $O(n)$ \\
    Ciur pentru $\tau$, $\sigma$      & $O(n \log n)$ & $O(n)$ \\
    \bottomrule
  \end{tabular}
  \end{center}
\end{frame}

\begin{frame}{Probleme de antrenament}
  \begin{enumerate}
    \item \href{https://code.algolymp.com/problems/4229}{\textbf{Prime and Divisor Count Queries}}
          \textcolor{gray}{(AlgoCode \#4229)}
    \item \href{https://code.algolymp.com/problems/737}{\textbf{Prime Factorization Exponents}}
          \textcolor{gray}{(AlgoCode \#737)}
    \item \href{https://code.algolymp.com/problems/727}{\textbf{Prime Factor with Largest Exponent}}
          \textcolor{gray}{(AlgoCode \#727)}
    \item \href{https://code.algolymp.com/problems/4201}{\textbf{Divisor Summary}}
          \textcolor{gray}{(AlgoCode \#4201)}
  \end{enumerate}
\end{frame}

\begin{frame}{Resurse}
  \begin{enumerate}
    \item \href{https://infoas.ro/lectie/91/complexitatea-unui-algoritm-timp-si-spatiu-in-c}{\textbf{infoas.ro}}
          \textemdash{} complexitate timp și spațiu în C++
    \item \href{https://www.bigocheatsheet.com/}{\textbf{bigocheatsheet.com}}
          \textemdash{} tabel vizual cu complexitățile comune
    \item \href{https://www.pbinfo.ro/articole/23702/standard-template-library-stl}{\textbf{pbinfo.ro}}
          \textemdash{} Standard Template Library
    \item \href{https://cppreference.com}{\textbf{cppreference.com}}
          \textemdash{} documentație completă C++
  \end{enumerate}
\end{frame}

\end{document}
